\subsection{Verbindungs- und Ortsvektoren}
\label{subsec:analytic_geometry_connecting_and_position_vectors}

Vektoren übersetzen die relative Lage zweier Punkte in eine Richtung mit
Länge. Mit einem festen Ursprung lässt sich außerdem jeder Punkt durch genau
einen Ortsvektor beschreiben.

\begin{DefinitionBox}{Verbindungsvektor}{analytic_geometry_connecting_vector}
  Der \emph{Verbindungsvektor} \EnglishTerm{displacement vector} vom Punkt
  $P_1(x_1,y_1,z_1)$ zum Punkt $P_2(x_2,y_2,z_2)$ ist
  \[
    \overrightarrow{P_1P_2}
    =\vec p_2-\vec p_1
    =\begin{pmatrix}
       x_2-x_1\\ y_2-y_1\\ z_2-z_1
     \end{pmatrix}.
  \]
  Sein Aufpunkt ist $P_1$, seine Pfeilspitze liegt in $P_2$.
\end{DefinitionBox}

\begin{KeyPointBox}{Abstand zweier Punkte}{analytic_geometry_point_distance}
  Der Abstand $d(P_1,P_2)$ ist die Länge des Verbindungsvektors:
  \[
    d(P_1,P_2)
    =\left|\overrightarrow{P_1P_2}\right|
    =\sqrt{(x_2-x_1)^2+(y_2-y_1)^2+(z_2-z_1)^2}.
  \]
\end{KeyPointBox}

\begin{DefinitionBox}{Ortsvektor}{analytic_geometry_position_vector}
  Der \emph{Ortsvektor} \EnglishTerm{position vector} eines Punktes
  $P(x,y,z)$ ist der Verbindungsvektor vom Ursprung $O(0,0,0)$ zu $P$:
  \[
    \vec p=\overrightarrow{OP}=\begin{pmatrix}x\\y\\z\end{pmatrix}.
  \]
  Sein Aufpunkt ist immer der Ursprung.
\end{DefinitionBox}

\begin{center}
  \CachedTikzFigure{content/resources/chapters/analytic_geometry/tikz/connecting_position_vectors.tex}
\end{center}

\begin{ExampleBox}{Verbindungsvektor und Punktabstand}{analytic_geometry_point_distance_example}
  Gegeben sind $P_1(1,1)$ und $P_2(3,3)$. Der Verbindungsvektor lautet
  \[
    \overrightarrow{P_1P_2}
    =\begin{ClozeContinuation}[11]{3}
      \begin{pmatrix}3-1\\3-1\end{pmatrix}
      =\begin{pmatrix}2\\2\end{pmatrix}
    \end{ClozeContinuation}.
  \]
  Damit ergibt sich der Abstand zu
  \[
    d(P_1,P_2)
    =\Cloze{7cm}{\sqrt{2^2+2^2}=\sqrt 8=2\sqrt 2}.
  \]
  Die Rechnung bestimmt also zugleich Richtung und Länge der Verbindung.
  \begin{center}
    \CachedTikzFigure{content/resources/chapters/analytic_geometry/tikz/point_distance_2d.tex}
  \end{center}
\end{ExampleBox}
